\section{Introdução}
\label{sec-introducao}

A crescente competitividade do setor de alimentos e bebidas exige que as empresas aprimorem continuamente seus processos produtivos para reduzir custos, aumentar a qualidade dos produtos e responder de forma ágil às oscilações do mercado. Nesse contexto, a compreensão dos fatores que influenciam a eficiência da produção tornam-se fundamental para o desenvolvimento de estratégias que conciliem competitividade, qualidade e rentabilidade \cite{rochaimportancia}. Entre essas estratégias, destaca-se a produção puxada, princípio da manufatura enxuta no qual a produção é orientada pela demanda do cliente, permitindo maior flexibilidade operacional e adaptação às variações do mercado \cite{tovsanovic2021evaluation}.

Em empresas produtoras de sucos naturais, a elevada variabilidade da demanda representa um desafio para o planejamento da produção. Como o sistema produtivo é baseado na produção puxada, imprecisões nas estimativas de vendas afetam diretamente o dimensionamento dos estoques, a aquisição de insumos e a utilização da capacidade produtiva. Além disso, a perecibilidade das matérias-primas torna o excesso de estoque custoso, enquanto previsões subestimadas podem ocasionar rupturas no abastecimento e perda de vendas \cite{dandaro2015planejamento}.

Nesse cenário, técnicas de Inteligência Artificial (IA) têm se consolidado como ferramentas promissoras para apoiar a tomada de decisão. A partir de dados históricos de vendas e de variáveis explicativas, como sazonalidade, calendário e fatores externos, modelos preditivos são capazes de identificar padrões de comportamento da demanda e gerar estimativas mais precisas para o planejamento operacional \cite{douaioui2024machine}. Além de reduzir a dependência de estimativas subjetivas, esses modelos podem ser continuamente atualizados com novos dados, adaptando-se às mudanças no comportamento dos consumidores.

Entre as abordagens utilizadas para previsão de demanda, a regressão linear destaca-se pela simplicidade, interpretabilidade e eficiência na modelagem de relações lineares entre variáveis \cite{montgomery2021introduction}. Em contrapartida, modelos de aprendizado profundo, como redes neurais recorrentes do tipo \textit{Long Short-Term Memory} (LSTM), apresentam maior capacidade para representar relações não lineares e dependências temporais complexas, características frequentemente presentes em séries históricas de vendas \cite{james2023statistical}. Dessa forma, a comparação entre essas abordagens possibilita avaliar o equilíbrio entre simplicidade, interpretabilidade e capacidade preditiva em diferentes contextos.

Diante desse contexto, este trabalho tem como objetivo desenvolver e comparar modelos de previsão de demanda para estimar o volume de vendas do mês subsequente em uma empresa de sucos naturais. Para isso, são utilizados dados históricos de vendas para a construção e avaliação de modelos baseados em regressão linear e redes neurais LSTM. Espera-se que os resultados obtidos contribuam para o planejamento da produção, reduzindo desperdícios, otimizando a alocação de recursos e fornecendo suporte à tomada de decisão baseada em dados.

\subsection{Trabalhos Relacionados}

A previsão de demanda por meio de aprendizado de máquina tem sido amplamente investigada, principalmente por meio da comparação entre modelos estatísticos tradicionais e arquiteturas de aprendizado profundo. Entre estas, as redes neurais LSTM destacam-se por sua capacidade de capturar dependências temporais de longo prazo e padrões sazonais complexos, frequentemente superando modelos lineares em cenários com elevada variabilidade temporal \cite{douaioui2024machine, veres2024application}. Entretanto, essa superioridade não é universal. \cite{veres2024application} ressaltam que o desempenho de modelos mais complexos depende das características dos dados e do contexto de aplicação. De forma semelhante, \cite{golabek2020demand}, ao compararem diferentes modelos para previsão de demanda de produtos de consumo rápido, observaram que, na categoria de bebidas, algoritmos mais simples, como Regressão Linear e \textit{Random Forest}, alcançaram desempenho equivalente ou superior ao das LSTMs.

Outro desafio recorrente na literatura refere-se à limitação dos dados históricos disponíveis. Em cenários caracterizados por séries temporais curtas ou intermitentes, modelos convencionais podem apresentar desempenho insatisfatório. Nesse contexto, \cite{rovzanec2103reframing} propuseram uma abordagem que separa a previsão da ocorrência da demanda da previsão de sua magnitude, além de destacarem a importância da utilização de métricas de avaliação adequadas para evitar interpretações equivocadas dos resultados.

Este trabalho insere-se nesse contexto ao investigar a previsão de demanda em uma pequena empresa do setor de bebidas naturais, caracterizada por histórico de vendas limitado e diferenças expressivas de volume entre os sabores comercializados. Especificamente, o estudo compara o desempenho da regressão linear e da LSTM em um cenário de baixa disponibilidade de dados, analisa a influência do volume de informações históricas na qualidade das previsões e discute uma limitação metodológica pouco explorada na literatura, relacionada à interpretação de métricas agregadas de desempenho quando aplicadas a produtos com escalas de vendas heterogêneas. Nesse caso, demonstra-se que métricas globais podem mascarar o desempenho preditivo em sabores individuais, evidenciando a necessidade de avaliações complementares por grupo de produtos.